\chapter{The twiddle constants, proved and not copied}

\textbf{Purpose:} State what it costs to prove a transform's constants instead of inheriting them, what that proof catches, and why every answer it gives carries the witness establishing it. \textbf{Scope:} Proth's theorem, the order of a root of unity, the group law on a twiddle table, and two silent wrong answers found by running it.

A number theoretic transform is a sum of powers of one constant. Everything the transform does depends on that constant having exactly the order it is supposed to have, and almost nobody checks.

\section{What a twiddle constant is}

The transform of a sequence of length \(n\) over the integers modulo \(q\) is
\[
\hat{p}_j = \sum_i p_i \, \omega^{ij}
\]
where \(\omega\) is an \(n\)-th root of unity modulo \(q\). The powers of \(\omega\) the butterflies evaluate are the \textbf{twiddle constants}. In the Cooley-Tukey butterfly they enter as \(c = a + bw\) and \(d = a - bw\); in the Gentleman-Sande butterfly as \(c = a + b\) and \(d = (a-b)w\). A transform of length \(2^k\) is \(k\) stages of \(n/2\) butterflies, so the twiddles are touched more often than any other value in the computation and are the only inputs supposed to be constant.

\section{The attack, and why it works}

Ravi, Yang, Bhasin, Zhang and Chattopadhyay inject a single electromagnetic fault setting the twiddle constants to zero. The transform does not fail. It returns a value of the ordinary shape carrying far less entropy than it should, and that is enough to recover a Kyber key and to forge a Dilithium signature, including a bypass of Dilithium's own verification.

The attack works because \textbf{nothing downstream can tell a real root of unity from a fiddled one}. The output of a transform with broken twiddles is the same size, the same type and the same range as the output of a correct one. There is no error to raise. A wrong twiddle is not detected, it is used.

\section{The same fault, with no attacker}

This work produced that fault condition by carelessness, which is worth recording because it is the cheaper way to arrive at it.

A transform was built on three moduli. The third, \(1610612737\), was taken on faith because it has the right shape, \(3 \cdot 2^{29} + 1\), the form such a modulus must have. Its generator was raised to \((p-1)/L\) to produce a root of unity of order \(L\), exactly as the construction requires, and a check written beside it reported the root was fine.

The check was trial division to one hundred, on a ten digit number. \(1610612737\) is \textbf{composite}, and the proof is one line: \(2^{1610612736} \not\equiv 1 \pmod{1610612737}\), so Fermat refuses it. The certificate is the witness \(2\).

Because it is composite, \(g^{(p-1)/L}\) had no particular order, the twiddles were wrong, and the transform returned numbers of the ordinary shape that were not the convolution of anything. The two correct moduli in the same set returned zero wrong coefficients out of thirty-two. The composite one returned thirty-two out of thirty-two, and the only reason anybody noticed is that the three were recombined against each other and disagreed.

\textbf{A copied constant is a constant nobody checked.} The moduli general purpose libraries reach for are the same values everywhere, carried between codebases as literals. The tables are correct, and that is not the point. The point is that correctness is being inherited instead of established, and inherited correctness cannot tell you when a value has been fiddled.

\section{Proth's theorem, which makes the check a proof}

\begin{theorem}[Proth]
Let \(N = k \cdot 2^n + 1\) with \(k\) odd and \(k < 2^n\). Then \(N\) is prime if and only if there is an integer \(a\) with \(a^{(N-1)/2} \equiv -1 \pmod N\).
\end{theorem}

Every modulus such a transform can use has this shape, because admitting a root of unity of order \(2^n\) is the same as \(2^n\) dividing \(N-1\). The theorem therefore reaches every candidate that matters.

\textbf{Why a witness is a proof.} Suppose \(a^{(N-1)/2} \equiv -1 \pmod N\), and let \(p\) be any prime factor of \(N\). Then \(a^{(N-1)/2} \equiv -1 \pmod p\) as well, so the order of \(a\) modulo \(p\) divides \(N-1\) and does not divide \((N-1)/2 = k \cdot 2^{n-1}\). Since \(N - 1 = k \cdot 2^n\), the order must carry the full \(2^n\), giving \(2^n \mid p-1\) and so \(p \equiv 1 \pmod{2^n}\). Every prime factor of \(N\) is at least \(2^n + 1\).

If \(N\) were composite it would have at least two such factors, so \(N \ge (2^n+1)^2 > 2^{2n}\). But \(k < 2^n\) gives \(N = k \cdot 2^n + 1 < 2^{2n} + 1\). The two cannot both hold, so \(N\) is prime. \(\square\)

The condition \(k < 2^n\) is what that last step needs, and a modulus outside that range is reported as out of range instead of being tested anyway.

\textbf{Why a witness is easy to find.} If \(N\) is prime then Euler's criterion makes \(a^{(N-1)/2}\) the Legendre symbol of \(a\), which is \(-1\) for every quadratic non-residue, and that is half of all \(a\). A witness is therefore found by trying small integers, and the first one usually works.

\textbf{Both verdicts are certificates.} A witness \(a\) with \(a^{(N-1)/2} \equiv -1\) proves primality. A witness \(a\) with \(a^{N-1} \not\equiv 1\) proves compositeness, by Fermat. Neither answer is a probability and neither is a judgement call. The only other outcome is running out of candidates, reported as inconclusive instead of rounded toward either verdict.

\section{The order of the root, exactly, without factoring anything}

Proving the modulus prime is not enough. The twiddle is \(\omega = g^{(p-1)/L}\), and the step from a generator to a root of unity is where the constant stops being checked.

\begin{proposition}
For \(L = 2^m\), if \(\omega^L = 1\) and \(\omega^{L/2} \ne 1\), then \(\omega\) has order exactly \(L\).
\end{proposition}

\textbf{Proof.} The order divides \(L = 2^m\), so it is \(2^j\) for some \(j \le m\). If \(j < m\) then \(2^j\) divides \(2^{m-1} = L/2\), which would give \(\omega^{L/2} = 1\). It does not, so \(j = m\). \(\square\)

Two exponentiations settle it completely, with no factoring of anything. This is cheap enough to run on every twiddle before every transform, and it is the step the fault attack removes. It refuses a zeroed twiddle, a root of short order, and a root drawn from a composite modulus, all by the same test.

The generator is proved the same way and just as completely: \(g\) generates the whole group if and only if \(g^{(p-1)/q} \ne 1\) for every prime \(q\) dividing \(p-1\). Here \(p - 1 = k \cdot 2^n\) with \(k\) small, so factoring it is factoring \(k\), and the proof is complete instead of partial.

\section{What it costs on one desktop}

One proof is one modular exponentiation, measured on unbounded integers with nothing vectorised and nothing on the card:

\begin{center}
\begin{tabular}{@{}rrr@{}}
\toprule
bits & decimal digits & one proof \\
\midrule
64 & 19 & \(0.0001\) s \\
128 & 39 & \(0.0001\) s \\
256 & 77 & \(0.0002\) s \\
512 & 154 & \(0.0011\) s \\
1024 & 308 & \(0.0185\) s \\
2048 & 617 & \(0.0419\) s \\
4096 & \(1{,}233\) & \(0.2876\) s \\
\bottomrule
\end{tabular}
\end{center}

A \textbf{1,233 digit number proved prime in under three tenths of a second}, deterministically. Not a probable prime, not a result with a confidence attached. A certificate, with the witness recorded beside it, letting a reader recheck the claim without rerunning anything.

The cost grows as roughly the square of the bit length, since the exponent has \(N\) bits and each squaring costs the multiply. This is the same theorem the large prime searches run, where the records are Proth primes of several million digits, and the difference between those runs and this table is only how long one is willing to wait.

\section{Why the size is ours to choose}

The tabled moduli are small because they are sized to a machine word. A lane holds sixty-four bits, a product of two residues has to fit in one, and the modulus therefore stops below \(2^{32}\). That is a property of the lane and not of the mathematics.

The arithmetic here is on unbounded integers, so the proof runs at any width, and moduli far past anything tabled can be produced and proved on demand:

\begin{center}
\begin{tabular}{@{}rlr@{}}
\toprule
modulus & shape & transform length \\
\midrule
\(18446744069414584321\) & \((2^{32}-1) \cdot 2^{32} + 1\) & \(2^{32}\) \\
\(9223372195768565761\) & \(2147483685 \cdot 2^{32} + 1\) & \(2^{32}\) \\
\(39614081257148775820397903873\) & \(140737488355387 \cdot 2^{48} + 1\) & \(2^{48}\) \\
\bottomrule
\end{tabular}
\end{center}

The last admits a transform of length \(281{,}474{,}976{,}710{,}656\). Nobody tables that, because it does not fit in a lane. \textbf{The floor belongs to the format, and the format is a choice.}

\section{The group law, and the proof that it is not enough on its own}

The twiddle table is not a list of numbers that happen to be useful. It is a group, and saying so precisely is what licenses the way the transform is implemented.

\begin{proposition}[The law]
Let \(\omega\) have order exactly \(n\) modulo a prime \(p\), and let \(w_j = \omega^j\) for \(0 \le j < n\). Then \(\{w_j\}\) is the cyclic group of order \(n\) generated by \(\omega\), and \(w_j w_k = w_{(j+k) \bmod n}\) for every \(j\) and \(k\).
\end{proposition}

\textbf{Proof.} \(w_j w_k = \omega^{j+k}\). Write \(j + k = qn + r\) with \(0 \le r < n\). Then \(\omega^{j+k} = (\omega^n)^q \omega^r = \omega^r = w_r\), and \(r\) is \((j+k) \bmod n\). The identity is \(w_0 = 1\) and the inverse of \(w_j\) is \(w_{n-j}\), so the set is closed, has an identity and has inverses, which makes it a group; it is generated by one element, so it is cyclic. \(\square\)

Two consequences are used directly.

\textbf{The table sums to zero.} For \(n > 1\), \(\sum_j w_j = (\omega^n - 1)/(\omega - 1) = 0\), exactly, because \(\omega \ne 1\) makes the denominator invertible. This is an integer zero and not a small residual.

\textbf{The table is generated by repeated multiplication.} \(w_{j+1} = w_j \omega\), so building the whole table costs one modular multiply per entry. In floating point that construction is warned against because error accumulates along the chain. In the exact ring there is no error to accumulate, so the cheap construction and the correct one are the same construction.

\subsection{The disproof: the law is necessary and is not sufficient}

\begin{proposition}
Let \(\omega'\) have order \(d\) where \(d\) divides \(n\) and \(d < n\). Then the table \(w'_j = \omega'^j\) satisfies the group law at length \(n\) exactly, while being the wrong table.
\end{proposition}

\textbf{Proof.} Since \(d\) divides \(n\), \(\omega'^n = (\omega'^d)^{n/d} = 1\). The argument above then runs unchanged with \(\omega'\) in place of \(\omega\), and \(w'_j w'_k = w'_{(j+k) \bmod n}\) holds for every \(j\) and \(k\). \(\square\)

A root of half the required order therefore produces a table that is a genuine homomorphic image of the wrong group. At length 4096 with \(\omega'\) of order 2048:

\begin{center}
\begin{tabular}{@{}ll@{}}
\toprule
invariant & verdict on a root of half the order \\
\midrule
\(w_0 = 1\) & passes \\
\(\sum w_j = 0\) & passes \\
\(\sum w_j^2 = 0\) & passes \\
group law on sampled pairs & passes, 0 of 256 failures \\
\textbf{order test} & \textbf{fails, the only one that does} \\
\bottomrule
\end{tabular}
\end{center}

Every invariant computable from the table at \(O(n)\) cost is blind to it, because each one is a statement the smaller group also satisfies. What separates them is the order, and the order is not a property any single relation among entries exposes. The separation costs two exponentiations, and it catches the case every cheaper test misses.

\subsection{Why the order of the checks matters operationally}

The transform builds its twiddles once by repeated multiplication and reads them at a stride, and the group law is what permits that. But the group law permits the construction from \emph{any} root whose order divides the length, including a wrong one, and the result would be a transform returning values of the ordinary shape that are not the convolution of anything.

The order proof is therefore not an optional extra beside the optimisation. It is the precondition for the optimisation being safe, and it has to run first. Stated as a rule: \textbf{prove the order, then generate the table.} Reversing those two produces a fast wrong answer.

The change was measured. Replacing an exponentiation per butterfly with one table read cut the transform at length \(2^{27}\) from \(2006\) ms to \(680\) ms, a factor of \(2.95\), and the output digest was identical at every length tested. A correct optimisation moves the clock and not the answer, and that check says which one happened.

\section{On the card, where a second silent wrong answer was waiting}

The device transform runs over three moduli, each carrying its Proth witness beside it in the source:

\begin{center}
\begin{tabular}{@{}rlrr@{}}
\toprule
modulus & shape & witness & order \\
\midrule
\(2013265921\) & \(15 \cdot 2^{27} + 1\) & 11 & \(2^{27}\) \\
\(2281701377\) & \(17 \cdot 2^{27} + 1\) & 3 & \(2^{27}\) \\
\(3892314113\) & \(29 \cdot 2^{27} + 1\) & 3 & \(2^{27}\) \\
\bottomrule
\end{tabular}
\end{center}

The first run of that kernel agreed with the host on the first modulus and differed in \emph{every coefficient} on the other two. The property separating them is not primality, which all three have. It is that \(2013265921\) is below \(2^{31}\) and the other two are above it. In the butterfly, two residues below a modulus above \(2^{31}\) sum past \(2^{32}\) and wrap, silently, in a type that raises nothing. The same is true of the borrow in the difference. Only the smallest modulus never reached the boundary, so one arm agreed and two did not.

\textbf{This is the same fault as the composite modulus, one level down.} In both cases an assumption about the arithmetic went unstated, the transform returned values of the ordinary shape, and nothing in the output said otherwise. The first was a constant nobody proved; the second was a width nobody checked. Neither was found by reading the code. Both were found because the gate ran cases differing in one property, so the failure named its own cause.

Widened to 64 bits, all three moduli agree with the host, and the device multiply agrees with the reference at every size tested.

\section{What it costs on the card}

One full multiply is two forward transforms, a pointwise product and one inverse, per modulus:

\begin{center}
\begin{tabular}{@{}rrr@{}}
\toprule
transform length & decimal digits carried & one modulus \\
\midrule
\(2^{20}\) & 10.1 million & \(8.5\) ms \\
\(2^{22}\) & 40.4 million & \(49.4\) ms \\
\(2^{24}\) & 162 million & \(220\) ms \\
\(2^{26}\) & 646 million & \(938\) ms \\
\(2^{27}\) & 1.29 billion & \(2006\) ms \\
\bottomrule
\end{tabular}
\end{center}

Three moduli at \(2^{27}\) is about six seconds of device time for a product of a billion decimal digits. The host's own multiply is Karatsuba, measured at \(0.930\) s for a million digits and \(5.610\) s for three million, an exponent of \(1.635\); extrapolated across three orders of magnitude that is some twenty hours for the same product. The extrapolation is stated as an extrapolation, since it reaches well past anything measured directly.

\textbf{The packing costs nothing, and that part is worth keeping.} A large integer is already a run of 32-bit limbs. Handing one to the device and reading one back are the same bytes named twice, and not a conversion. The device returns the product as three arrays and the caller assembles them at a 32-bit stride. Convolution coefficients run past \(2^{32}\) and have to be carried, and the carrying happens inside that addition, in the host's own big-integer arithmetic. A carry pass written out over a hundred million coefficients would cost more than every transform put together; here there is no pass at all.

\section{What this refuses}

A modulus outside Proth's range, where \(k\) is not below \(2^n\), since the converse is what a witness rests on and it does not hold there. A root whose order falls short of the length asked for, the fiddled twiddle exactly. And any verdict without its certificate: every answer carries the witness establishing it.

\section{Source}

Prasanna Ravi, Bolin Yang, Shivam Bhasin, Fan Zhang, Anupam Chattopadhyay. \emph{Fiddling the Twiddle Constants: Fault Injection Analysis of the Number Theoretic Transform.} IACR ePrint 2022/824.
